\subsubsection*{3.7.}
当 $\operatorname{char}(k)\ne2$ 时, 以及当 $X$ 的维数为偶数时, 定理 2.5 成立.

\noindent\underline{证明.}
鉴于 3.5.0 及 3.3, 只需验证以下断言:

\medskip
\noindent(3.7.1) 对每个闭点 $x_0\in X$ 和每个整数 $d\geq2$, 存在一个 $d$ 次超曲面 $H$, 它在 $x_0$ 处与 $X$ 相切, 且 $x_0$ 是 $X\cap H$ 的普通二次奇点.

选择齐次坐标 $X_0,\ldots,X_r$, 使 $x_0=(1,0,\ldots,0)$, 且函数 $X_i/X_0$, $i=1,\ldots,n$, 为 $X$ 在 $x_0$ 附近的局部坐标. 可取 $H$ 为由方程
\begin{equation}\tag{3.7.1.1}
H=X_0^{d-2}\sum_{i=1}^{n}X_i^2
\qquad\text{若 }\operatorname{char}(k)\ne2,
\end{equation}
或
\begin{equation}\tag{3.7.1.2}
H=X_0^{d-2}\sum_{i=1}^{n/2}X_iX_{n/2+i}
\qquad\text{若 }n\text{ 为偶数}.
\end{equation}
定义的超曲面. 当 $n$ 为奇数时(在第一种假设下), 源文记录的等价标准选择为
\begin{equation}\tag{3.7.1.3}
H=X_0^{d-2}
\left(X_n^2+
\sum_{i=1}^{\lfloor n/2\rfloor}
X_iX_{\lfloor n/2\rfloor+i}\right).
\end{equation}

对断言 2.5.2, 注意或者 $\varphi$ 一般非分歧, 此时可应用 3.5.0; 或者它处处分歧. 在特征零下, 后一种情形蕴含 $\dim\check X\leq r-2$, 故判别条件 (3.2.4) 自动成立.\footnote{印本此处写为 $\dim X\leq r-2$. 推论 3.5.0 中的另一种情形以及所引用的判别条件要求应为 $\dim\check X\leq r-2$.}